% Источники: exam/materials/моделирование.pdf, раздел 25; https://github.com/HumsterProgrammer/iu9-notes/blob/main/8sem/Моделирование/Моделирование_лекция-05_26-02-26.md.
\section{Организация работы управляющей программы моделирования (УПМ).}

\textbf{Управляющая программа моделирования (УПМ)} -- компонент имитационной модели, который координирует выполнение всех процессов, обеспечивает синхронизацию, управляет временем моделирования и собирает статистику.

Функции, реализуемые УПМ в специализированных средах моделирования:

\begin{enumerate}
  \item внутренняя синхронизация компонент модели;
  \item внешняя синхронизация компонент модели;
  \item уточнение исходной информации при моделировании;
  \item контроль за ходом имитации;
  \item организация окончания имитации;
  \item организация сбора статистики;
  \item документирование.
\end{enumerate}

\textbf{Уточнение исходной информации при моделировании.} Разрабатывается программный блок УПМ для задания начальных условий имитации. Обращение к блоку происходит однократно в модельный момент времени $t_0$.

\textbf{Контроль за ходом имитации.} Реализуется соответствующий блок; в оперативную память копируется текущее состояние модели, работа с которой может приостанавливаться и возобновляться.

\textbf{Организация окончания имитации.} Организуется блок окончания моделирования: последовательность действий по контролю момента окончания счета, начала обработки результатов и планирования вычислительного эксперимента.

\textbf{Организация сбора статистики.} Устанавливаются наблюдатели модели -- операторы, фиксирующие выборочные значения, например среднее время жизни заявки в системе, среднее время обработки заявки, количество потерянных заявок. Замер осуществляется при смене кванта модельного времени.

\textbf{Документирование.} Описание имитационной модели должно быть зафиксировано в отчете в общепринятом виде.

\clearpage
